\chapter{End-to-End Project: The GitHub Code Review Agent}

\begin{quote}
\textit{``The true test of a system is not whether it works in a demo — it is whether it holds together at 3 AM on a Friday when the oncall engineer is asleep, the CI pipeline has a flaky test, and four developers have pushed PRs simultaneously. Engineering is about systems that work when you are not watching.''}
\end{quote}

This final chapter builds a complete, production-hardened agent from scratch. We will not shortcut anything. By the end, you will have a deployable GitHub Code Review Agent that:
\begin{enumerate}
    \item Receives GitHub Pull Request webhooks via a FastAPI server
    \item Parses changed files and uses an AST-aware chunker to identify semantically meaningful code units
    \item Queries a vector database to find \textit{similar past bugs} from the project's history
    \item Runs a specialized reviewer LLM that produces structured review comments
    \item Posts comments directly to the GitHub PR using the GitHub API
    \item Enforces a cost cap of \$0.05 per review and a latency SLO of P95 $< 30$ seconds
\end{enumerate}

\section{Architecture Overview}

\begin{figure}[ht]
\centering
\begin{tikzpicture}[
    scale=0.78, transform shape,
    node distance=0.9cm and 1.4cm,
    box/.style={draw, rectangle, rounded corners, minimum height=0.7cm,
      minimum width=2.5cm, align=center, fill=blue!8, font=\small},
    dbbox/.style={draw, cylinder, shape border rotate=90, minimum height=0.9cm,
      minimum width=2.2cm, align=center, fill=green!10, font=\small},
    extbox/.style={draw, rectangle, rounded corners, minimum height=0.7cm,
      minimum width=2.5cm, align=center, fill=gray!15, font=\small},
    arrow/.style={thick,->,>=stealth}
]
% External systems
\node[extbox] (gh_pr) at (-5, 0) {GitHub PR\\(Webhook)};
\node[extbox] (gh_api) at (-5, -5.5) {GitHub API\\(Post Comments)};

% Core pipeline
\node[box, fill=purple!12] (webhook) at (0, 0) {FastAPI Webhook\\Handler};
\node[box] (diff) at (0, -1.3) {Git Diff Parser\\+ AST Chunker};
\node[box, fill=yellow!12] (retriever) at (0, -2.6) {Vector Retriever\\(Similar Bugs)};
\node[box, fill=orange!12] (reviewer) at (0, -3.9) {Reviewer LLM\\(GPT-4o-mini)};
\node[box] (formatter) at (0, -5.2) {Comment Formatter\\+ GitHub Poster};

% Vector DB
\node[dbbox] (vdb) at (3.5, -2.6) {Chroma DB\\(Bug History)};

% Cost guard
\node[box, fill=red!10] (guard) at (3.5, -3.9) {Cost Guard\\(\$0.05 cap)};

% Arrows
\draw[arrow] (gh_pr) -- node[above, font=\tiny] {POST /webhook} (webhook);
\draw[arrow] (webhook) -- (diff);
\draw[arrow] (diff) -- (retriever);
\draw[arrow] (retriever) -- (reviewer);
\draw[arrow] (reviewer) -- (formatter);
\draw[arrow] (formatter) -- node[right, font=\tiny] {POST review} (gh_api);
\draw[arrow, dashed] (retriever) -- node[above, font=\tiny] {embed + search} (vdb);
\draw[arrow, dashed] (vdb) -- (retriever);
\draw[arrow, dashed] (guard) -- node[right, font=\tiny] {check} (reviewer);

\end{tikzpicture}
\caption{GitHub Code Review Agent architecture. The dashed lines show side-channel operations (vector search and cost checking) that do not block the main data flow.}
\end{figure}

\section{Step 1: The Webhook Server}


\begin{center}
\noindent\rule{0.6\textwidth}{0.4pt} \\
\vspace{0.3em}
\textit{[Code implementation is available in the companion coding handbook: \texttt{coding-handbook/ch18\_github\_agent/}]} \\
\noindent\rule{0.6\textwidth}{0.4pt}
\end{center}


\section{Step 2: AST-Aware Code Chunking}

Naive line-based chunking splits functions in half. AST-aware chunking respects Python syntax boundaries, producing semantically meaningful review units.


\begin{center}
\noindent\rule{0.6\textwidth}{0.4pt} \\
\vspace{0.3em}
\textit{[Code implementation is available in the companion coding handbook: \texttt{coding-handbook/ch18\_github\_agent/}]} \\
\noindent\rule{0.6\textwidth}{0.4pt}
\end{center}


\section{Step 3: Semantic Bug Retrieval from History}


\begin{center}
\noindent\rule{0.6\textwidth}{0.4pt} \\
\vspace{0.3em}
\textit{[Code implementation is available in the companion coding handbook: \texttt{coding-handbook/ch18\_github\_agent/}]} \\
\noindent\rule{0.6\textwidth}{0.4pt}
\end{center}


\section{Step 4: The Reviewer LLM with Structured Output}


\begin{center}
\noindent\rule{0.6\textwidth}{0.4pt} \\
\vspace{0.3em}
\textit{[Code implementation is available in the companion coding handbook: \texttt{coding-handbook/ch18\_github\_agent/}]} \\
\noindent\rule{0.6\textwidth}{0.4pt}
\end{center}


\section{Step 5: Posting to GitHub \& Full Pipeline Assembly}


\begin{center}
\noindent\rule{0.6\textwidth}{0.4pt} \\
\vspace{0.3em}
\textit{[Code implementation is available in the companion coding handbook: \texttt{coding-handbook/ch18\_github\_agent/}]} \\
\noindent\rule{0.6\textwidth}{0.4pt}
\end{center}


\section{Production Hardening Checklist}

Before deploying this agent to handle real production PRs, verify:

\begin{center}
\begin{tabular}{|l|c|l|}
\hline
\textbf{Requirement} & \textbf{Implemented?} & \textbf{Where} \\
\hline
Webhook signature verification & \ding{51} & \texttt{verify\_github\_signature()} \\
Async background processing (no timeout) & \ding{51} & \texttt{BackgroundTasks} \\
Cost cap per review & \ding{51} & \texttt{cost\_tracker} \\
Max chunk limit & \ding{51} & \texttt{chunks[:10]} \\
AST-aware chunking (no split functions) & \ding{51} & \texttt{ast\_chunk\_python\_file()} \\
Idempotent review (no double-posting) & \ding{55} & Add SHA-based deduplication \\
Observability (traces per PR) & \ding{55} & Add OpenTelemetry (Ch. 15) \\
Rate limit handling (GitHub 5000 req/hr) & \ding{55} & Add \texttt{with\_backoff()} (Ch. 3) \\
Secrets in environment (not hardcoded) & \ding{51} & \texttt{os.environ} \\
\hline
\end{tabular}
\end{center}

\textbf{The two unchecked items} — idempotency and observability — are your homework. The patterns for both are fully covered in Chapter 3 (SHA hashing of action signatures) and Chapter 15 (OpenTelemetry tracing). A production-grade agent is an integration of all of the preceding chapters.
